% エントリ 15(2026-11-10): 三準位の外の準位(411、433)と印 4
\section{[E15, 2026-11-10] 三準位の外の準位: 411・433 が育つ場所と、印 4 の付け直し}\label{sec:e15}

E14 までの印 4 は「$\alpha>0.15$ の固定点」を一律に指していた(ユーザー指示「433 を入れる」)。BGLS の Sec.~IV C を読み直すと、433 だけを入れても足りないことがわかった。
\begin{itemize}
\item \textbf{433}: (211, 322) の固定点の上で育つのは、$\alpha\tilde r_+^{1/3}\gtrsim0.31$ のときだけ(BGLS の n33 の式。$\kappa^{211\times\infty}_{322\times433}/\kappa^{433\times\mathrm{BH}}_{211\times322}=2.6\times10^{-9}/9.1\times10^{-8}$ の 3 乗根)。主要項の率では $\alpha\gtrsim0.25$。
\item \textbf{411}: BGLS は $\alpha\gtrsim0.2$ で育ち、211 の準平衡の値の最大 50 倍になると書く。そのあと $411\times411\to422\times\mathrm{BH}$(共鳴)で 422 が育つが、この率は表になく、BGLS も今後の課題としている。
\end{itemize}
そこで、固定点で 411・433・544 の 1 粒子あたりの増加率を計算し、それが BH の降下のあいだに 1 個の量子から育つだけの $e$-fold を稼ぐ場所だけを印 4 にした。計算は \texttt{calc/fp3.py}(増加率と五準位の ODE)、\texttt{calc/chisr.py}(\texttt{extra\_growth}、\texttt{flag\_fp})、\texttt{calc/e15\_extra\_levels.py}、出力は \texttt{calc/out/e15.json}。

\subsection{増加率}

固定点 $(\varepsilon_{211},\varepsilon_{322},\varepsilon_{544})$ で、BGLS の Table V・VI の率から
\begin{align}
\lambda_{411}&=\gamma_{411}-(\kappa^{322\times\mathrm{BH}}_{211\times411}+\kappa^{211\times\infty}_{322\times411})\varepsilon_{211}\varepsilon_{322},\\
\lambda_{433}&=\gamma_{433}+(\kappa^{433\times\mathrm{BH}}_{211\times322}-\kappa^{211\times\infty}_{322\times433})\varepsilon_{211}\varepsilon_{322}
-\kappa^{544\times\mathrm{BH}}_{211\times433}\varepsilon_{211}\varepsilon_{544}-(\kappa^{322\times\infty}_{433\times544}\varepsilon_{322}+\kappa^{211\times\infty}_{433\times544}\varepsilon_{211})\varepsilon_{544}
\end{align}
($\kappa$ には $(\Mpl/f)^4\mu$ と、BH が吸収する過程では $\alpha^{11}\tilde r_+$、放出では $\alpha^8$ を掛ける)。固定点の占有数は $(f/\Mpl)^2$ に比例するので、$\lambda$ は $f$ によらない。544 は 211 の臨界より上で (211, 322) の固定点に対する増加率を使う。

\paragraph{411 の超放射の率}
411 の相対論的な率の表はない。そこで Leaver の連分数法(E4)で 411 を直接解いた。相対論的な率と解析式の比は、211 と 411 でほぼ同じだった($\alpha=0.1$、0.2、0.25、0.3 で、211 が 0.95、1.49、1.81、3.04、411 が 0.96、1.54、1.92、3.40)。比べるべきは 211 との比なので、411 の解析式に 211 の補正の比を掛けて使う(違いは 12\% 以内)。

\paragraph{有限 $\alpha$ の補正}
主要項の率では、411 は $\alpha\le0.30$ のどこでも育たない。BGLS の「411 は $\alpha\gtrsim0.2$ で育つ」は、BH が吸収する過程の有限 $\alpha$ の補正(BGLS の数値では $\alpha\sim0.2$ で数倍)を入れると再現した。BGLS によれば、この補正は地平線の近くで決まるので、BH が吸収する過程どうしの比はあまり変えない。そこで E7 の $C_\mathrm{BH}$ を BH が吸収する過程すべてに掛けて評価した。

\begin{result}
増加率が正になる場所($\mu$ によらず $\alpha$ と $\chi$ で決まる):
\begin{itemize}
\item \textbf{主要項}: 211 の臨界より上で、433 だけが $\alpha\ge0.28$ で育つ。411 と 544 は $\alpha\le0.30$ で育たない。
\item \textbf{$C_\mathrm{BH}$ 込み}: 211 の臨界より上で、433 と 544 が $\alpha\gtrsim0.17$、411 が $\alpha\gtrsim0.20$ で育つ(211 の臨界のごく近くでは 411 は育たない)。BGLS の 411 と 433 の記述と合う。
\item \textbf{211 の臨界より下の (322, 544) の固定点では、どちらでも何も育たない}($\alpha\le0.30$)。したがって E12・E14 の 211 の臨界より下への降下と、322 から育つ BH は、$\alpha\le0.30$ まで三準位で閉じている。
\end{itemize}
\end{result}

\subsection{五準位の ODE での確かめ}

211・322・411・433・544 の五準位の ODE(上の過程すべてと $433\times433\to211\times\infty$、$411\times411\to322\times\mathrm{BH}$、1 個の量子から)を、BGLS の Fig.~(411 の例)と同じ $10\Msun$、$\alpha=0.22$、$\chi_0=0.9$ で解いた($f=10^{14}\GeV$)。

\begin{center}\small
\begin{tabular}{lccccc}
\toprule
年齢 [年] & 1 & 10 & 100 & 300 & 1000\\
\midrule
五準位、$C_\mathrm{BH}$ 込み & 0.8997 & 0.8885 & 0.8023 & 0.7495 & 0.7386\\
三準位、$C_\mathrm{BH}$ 込み & 0.8999 & 0.8992 & 0.8922 & 0.8772 & 0.8328\\
五準位、主要項 & 0.8997 & 0.8965 & 0.8667 & 0.8158 & 0.7530\\
$\chi_\mathrm{SR}$(三準位、主要項) & 0.8996 & 0.8964 & 0.8666 & 0.8158 & 0.7531\\
\bottomrule
\end{tabular}
\end{center}

\begin{result}
\begin{itemize}
\item 主要項では 411・433 は育たず、五準位の ODE は $\chi_\mathrm{SR}$ と $10^{-4}$ で合う。判定のとおりである。
\item $C_\mathrm{BH}$ 込みでは 411 が 1 年ほどで育ち、211 の約 12 倍になる(BGLS の「最大 50 倍」と同じ傾向)。433・544 も育つ。三準位だけなら $C_\mathrm{BH}$ のぶん回転を失うのが遅くなるが、411 などがそれを打ち消し、主要項の三準位($\chi_\mathrm{SR}$ の既定)より少し速い。この例では、印 4 の点の $\chi_\mathrm{SR}$ は最大 0.066 ずれる($10^2$--$3\times10^2$ 年)。
\end{itemize}
\end{result}

\subsection{新しい印 4}

固定点の領域で、三準位の外の準位が降下のあいだに稼ぐ $e$-fold
\begin{equation}
N=\int\lambda\,\mathrm{d}t=\Bigl(\frac{\Mpl}{f}\Bigr)^2\int_{\chi_T}^{\chi_0}\frac{\max(\lambda/\mu,0)}{|H|}\,\mathrm{d}\chi
\end{equation}
が $\ln(GM^2)\approx180$ を超えるとき印 4 とする($\lambda$ は 411・433・544 の最大、$C_\mathrm{BH}$ 込み)。積分は $f$ によらないので、$I$ と同じ格子で表にした。322 の臨界まで下がった BH は、そのあと 433 の超放射だけで育つ $e$-fold($\mu\gamma_{433}T$)も見る(控えめに年齢全体)。値は三準位のまま返す。

\begin{result}
\begin{itemize}
\item $\mu=0.71\times10^{-12}\eV$、$f=10^{12}\GeV$、6 Gyr では、211 から育つ $30$--$40\Msun$($\chi_0=0.8$)が印 4、GW241011 の主星、$25\Msun$、322 から育つ $45\Msun$($\chi_0=0.7$)、$50\Msun$(0.6)は印 1。
\item Beta(0.69, 2.28) の 1 世代目($8$--$50\Msun$、6 Gyr)では、印 4 は $f=10^{12}$、$10^{14}$、$10^{16}\GeV$ で 2.3、7.4、10.4\%(E14 では一律に 11.2\%)。$f$ が大きいほど、322 の臨界まで下がって 433 の超放射が効く BH が増える。
\item 単体版も同じ判定にし、20 万個の標本で値も印もプロジェクト版と一致した。
\end{itemize}
\end{result}

\begin{caveat}
\begin{itemize}
\item \textbf{印 4 の中身}: 411 が育ったあとの 422 の成長(BGLS が予言、率は表にない)と、433 の超放射による $m=3$ の回転の減少は入れていない。印 4 の値はこの例で 0.07 ほどずれうる。直すには Witte \& Mummery の $n\le5$ のコードとの照合が要る。
\item \textbf{判定は $C_\mathrm{BH}$ に依存する}: $C_\mathrm{BH}$ は 211$\times$211$\to$322$\times$BH の値(E7、図から読んだもの)で、ほかの過程にも同じ値を掛けた。主要項だけで判定すると印 4 はもっと少ない。
\item \textbf{411 の率}: 解析式に 211 の補正の比を掛けたもの(Leaver との違いは 12\% 以内)。表は作っていない。
\end{itemize}
\end{caveat}
